\documentclass[11pt]{article}

% ---------------------------------------------------------------------------
% Kotoba: A Content-Addressed Distributed Datalog Database
% arXiv-style preprint.  Build:  latexmk -pdf kotoba.tex
% ---------------------------------------------------------------------------

\usepackage[utf8]{inputenc}
\usepackage[T1]{fontenc}
\usepackage{lmodern}
\usepackage[margin=1in]{geometry}
\usepackage{amsmath,amssymb,amsthm}
\usepackage{mathtools}
\usepackage{booktabs}
\usepackage{array}
\usepackage{graphicx}
\usepackage{listings}
\usepackage{xcolor}
\usepackage{enumitem}
\usepackage{microtype}
\usepackage{authblk}
\usepackage[hidelinks]{hyperref}
\usepackage{cleveref}

% ---- listings style -------------------------------------------------------
\definecolor{codebg}{rgb}{0.97,0.97,0.97}
\lstset{
  basicstyle=\ttfamily\small,
  backgroundcolor=\color{codebg},
  breaklines=true,
  columns=fullflexible,
  keepspaces=true,
  frame=single,
  framesep=4pt,
  rulecolor=\color{gray!40},
  showstringspaces=false,
  numbers=none,
}

% ---- theorem environments -------------------------------------------------
\newtheorem{definition}{Definition}
\newtheorem{property}{Property}

\newcommand{\kotoba}{\textsc{Kotoba}}
\newcommand{\cid}{\textsf{CID}}
\newcommand{\eq}{\;\triangleq\;}

% ---------------------------------------------------------------------------
\title{\textbf{\kotoba: A Content-Addressed Distributed Datalog Database\\
for Decentralized AI Agent Systems}}

\author[1]{Jun Kawasaki}
\affil[1]{com-junkawasaki / etzhayyim}
\date{June 2026}

% ---------------------------------------------------------------------------
\begin{document}
\maketitle

\begin{abstract}
We present \kotoba, a distributed, content-addressed knowledge-graph database
designed for decentralized AI-agent systems. \kotoba\ unifies eight design
primitives that are usually provided by separate systems: (i) immutable,
Datomic-style 5-tuple datoms keyed by IPFS content identifiers
(\cid{}s); (ii) a four-index in-memory \emph{arrangement} (EAVT/AEVT/AVET/VAET)
mirrored by a history-independent \emph{Prolly Tree} on disk; (iii) incremental,
semi-naive Datalog evaluation over change deltas, exposed additionally through a
SPARQL~1.1 surface; (iv) Pregel bulk-synchronous-parallel (BSP) graph
computation over the datom graph; (v) capability-based authorization via
chain-agnostic CACAO (CAIP-74) delegation chains; (vi) end-to-end encryption of
agent messages using the Signal protocol (X3DH + Double Ratchet); (vii) sandboxed,
polyglot user logic via the WebAssembly Component Model, including a native
Kotoba/EDN-to-WASM compiler; and (viii) a tiered, optionally AES-256-GCM-sealed
block store backed by IPFS. We describe the formal data model, the storage and
indexing layers, the query and computation engines, the security architecture,
and the runtime. We report micro- and macro-benchmarks on commodity hardware
(an Apple~M4 laptop), including a sustained ingest rate of $5{,}222$
entities/second and CACAO-gated, replay-protected query throughput of
$12{,}753$ queries/second. We close with a five-axis analysis of the open
problems---authority, ordering, availability, conflict, and transport---that
shape the path from a single-writer deployment to a fully multi-writer
decentralized one.
\end{abstract}

\noindent\textbf{Keywords:} content addressing, Datalog, Datomic, IPFS, Prolly
Tree, Pregel, CACAO, Signal protocol, WebAssembly, decentralized systems.

% ===========================================================================
\section{Introduction}
\label{sec:intro}

Decentralized AI-agent systems place an unusual combination of demands on their
backing store. Agents must (a) read and write a shared, evolving knowledge graph
without a central coordinator; (b) reason over that graph with recursive,
rule-based queries; (c) authenticate and authorize one another using portable,
cross-chain credentials rather than server-issued sessions; (d) exchange
private messages with end-to-end confidentiality; and (e) execute untrusted,
user-supplied logic safely. No single existing system covers this surface:
graph databases lack content addressing and capability security; content-addressed
stores (IPFS) lack a query engine; capability frameworks lack a database; and
agent frameworks lack durable, verifiable persistence.

\kotoba\ is our attempt to provide one coherent substrate for all of these
concerns. Its design is summarized by the following informal equation, which
also serves as a table of contents for the system:
%
\begin{equation}
\label{eq:kotoba}
\kotoba \eq
\underbrace{\mathrm{Datom}[\cid/T]}_{\text{\cref{sec:datom}}}
\times
\underbrace{\mathrm{EAVT}}_{\text{\cref{sec:index}}}
\times
\underbrace{\mathrm{Pregel}[\mathrm{BSP}]}_{\text{\cref{sec:pregel}}}
\times
\underbrace{\mathrm{Datalog}[\Delta]}_{\text{\cref{sec:datalog}}}
\times
\underbrace{\mathrm{CACAO}}_{\text{\cref{sec:cacao}}}
\times
\mathrm{AT}
\times
\mathrm{LLM}
\times
\underbrace{\mathrm{WASM/WIT}}_{\text{\cref{sec:wasm}}}.
\end{equation}

\paragraph{Contributions.}
\begin{enumerate}[leftmargin=1.4em,itemsep=2pt]
  \item A \emph{content-addressed datom model} that gives every fact, index node,
    transaction, and program a deterministic IPFS \cid{}, making the database a
    Merkle-CRDT whose commit identity is replica-independent
    (\cref{sec:storage,sec:datom}).
  \item A dual hot/cold index: a four-arrangement in-memory cache for
    sub-microsecond lookups, mirrored by a history-independent Prolly Tree that
    is DAG-CBOR/IPLD-encoded and stored as IPFS blocks (\cref{sec:index}).
  \item An incremental Datalog engine using semi-naive fixpoint over deltas,
    surfaced both natively and through SPARQL~1.1, and lifted to distributed
    Pregel BSP computation (\cref{sec:datalog,sec:pregel}).
  \item A security architecture combining CACAO (CAIP-74) capability delegation,
    Signal end-to-end encryption, and an optionally AES-256-GCM-sealed cold tier
    with content-derived nonces (\cref{sec:security}).
  \item A WebAssembly Component-Model runtime with a typed WIT host interface and
    a native Kotoba/EDN-to-WASM compiler, enabling sandboxed polyglot user-defined
    functions (\cref{sec:wasm}).
  \item An empirical evaluation on commodity hardware and a five-axis roadmap
    toward multi-writer decentralization (\cref{sec:eval,sec:roadmap}).
\end{enumerate}

% ===========================================================================
\section{System Overview}
\label{sec:overview}

\kotoba\ is implemented in Rust as a workspace of roughly thirty crates. The
layering follows the data path: a content-addressing and storage core
(\texttt{kotoba-core}, \texttt{kotoba-store}); a graph and datom layer
(\texttt{kotoba-graph}, \texttt{kotoba-datomic}); query and computation
(\texttt{kotoba-query}, \texttt{kotoba-vm}); security
(\texttt{kotoba-auth}, \texttt{kotoba-signal}, \texttt{kotoba-crypto},
\texttt{kotoba-custody}); a WASM runtime and compilers
(\texttt{kotoba-runtime}, \texttt{kotoba-clj}, \texttt{kotoba-guest}); and
networking and service surfaces (\texttt{kotoba-net}, \texttt{kotoba-dht},
\texttt{kotoba-server}). \Cref{tab:crates} lists the principal crates and their
responsibilities.

\begin{table}[t]
\centering
\small
\caption{Principal \kotoba\ crates and their responsibilities.}
\label{tab:crates}
\begin{tabular}{@{}p{0.22\linewidth} p{0.70\linewidth}@{}}
\toprule
\textbf{Crate} & \textbf{Responsibility} \\
\midrule
\texttt{kotoba-core}    & CIDv1, frame encoding, Prolly Tree storage. \\
\texttt{kotoba-store}   & Block stores: in-memory, Kubo/IPFS HTTP, tiered, sealed (AES-256-GCM). \\
\texttt{kotoba-graph}   & Quad/datom API, CommitDag, SPARQL 1.1, $N$-hop \texttt{DESCRIBE}. \\
\texttt{kotoba-datomic} & Datomic-compatible schema and transact semantics. \\
\texttt{kotoba-query}   & Four-index arrangement, incremental Datalog, materialized views. \\
\texttt{kotoba-vm}      & Pregel BSP, state graph, Datalog executor. \\
\texttt{kotoba-auth}    & CACAO chain verification; EdDSA, EIP-191/1271, BIP-322. \\
\texttt{kotoba-signal}  & Signal protocol: X3DH key agreement, Double Ratchet. \\
\texttt{kotoba-custody} & $t$-of-$N$ Shamir custody of the block-sealing key. \\
\texttt{kotoba-runtime} & WASM Component-Model host and WIT bindings. \\
\texttt{kotoba-clj}     & Kotoba/EDN $\rightarrow$ WebAssembly compiler. \\
\texttt{kotoba-net}, \texttt{kotoba-dht} & libp2p QUIC/Noise/GossipSub, Kademlia DHT, firehose. \\
\texttt{kotoba-server}  & XRPC/MCP HTTP endpoints, streaming firehose. \\
\bottomrule
\end{tabular}
\end{table}

% ===========================================================================
\section{Content-Addressed Storage}
\label{sec:storage}

\subsection{Content identifiers}
Every byte string persisted by \kotoba\ is named by a \cid: a self-describing
multihash. We use CIDv1 with the \texttt{sha2-256} hash function and the
\texttt{dag-cbor} codec, making identifiers byte-compatible with IPFS and IPLD
links (tag-42). Because the identifier is the hash of the canonical encoding,
naming is deterministic and replica-independent: two nodes that compute the same
logical object derive the same \cid\ without any coordination.

\subsection{Prolly Tree}
\label{sec:prolly}
The on-disk index is a \emph{Prolly Tree} (a probabilistic, content-defined
Merkle B-tree). Node boundaries are placed by content rather than position: a
key $k$ ends a chunk when
\[
\mathrm{blake3}(k)[0..4]\;\&\;\mathtt{BOUNDARY\_MASK} = 0,
\]
where $\mathtt{BOUNDARY\_MASK}=\mathtt{0xFF}$ gives a boundary probability of
$\approx 1/256$ and hence chunks of $\approx 256$ entries on average (the mask is
tunable). Leaf and internal nodes are
%
\begin{lstlisting}
enum ProllyNode {
  Leaf     { entries:  Vec<(Vec<u8>, Vec<u8>)>,    cid: KotobaCid },
  Internal { children: Vec<(Vec<u8>, KotobaCid)>,  cid: KotobaCid },
}
\end{lstlisting}
%
each canonically DAG-CBOR-encoded and addressed by its \cid. This construction
gives two properties we rely on throughout the system.

\begin{property}[History independence]
The tree root \cid\ is a function of the entry \emph{set} only, independent of
insertion order. Two replicas that ingest the same facts in any order converge
to the same root.
\end{property}

\begin{property}[Structural sharing]
A delta commit rewrites only the nodes on the affected root-to-leaf paths;
unchanged subtrees retain their \cid s and are not re-encoded. We exploit this to
build the four indices in parallel against a capturing block store, which yields
roughly a $28\%$ ingest speedup over a serial rebuild.
\end{property}

\subsection{Tiered and sealed block store}
\label{sec:store}
The block store is a composable hierarchy. \texttt{MemoryBlockStore} serves
microsecond hot lookups; \texttt{KuboBlockStore} talks HTTP to a Kubo/IPFS node
for durable cold storage; \texttt{BudgetedBlockStore} caps the hot tier with an
LRU budget; and \texttt{TieredBlockStore} composes a hot and a cold store so that
writes go to both (cold asynchronously) and reads fall through, promoting cold
hits back to the hot tier. On \texttt{wasm32}, \texttt{IdbBlockStore} backs the
store with browser IndexedDB.

For confidentiality at rest, \texttt{SealedBlockStore} wraps any cold store and
encrypts every exported block with AES-256-GCM under a node-held key. The nonce
is \emph{content-derived}, $\mathrm{nonce}=\mathrm{HKDF}(\text{block key},
\text{``}\dots\text{nonce/v1''}, \text{plaintext }\cid)$, which is deterministic
yet collision-free, so re-sealing identical content is idempotent and no nonce is
ever reused. Sealed blocks carry a \texttt{KSB1} magic prefix and are themselves
re-addressed by the \cid\ of their ciphertext, so that only ciphertext is pinned
to IPFS and fanned out to peers. A sidecar index of (plaintext \cid, sealed \cid)
pairs supports deterministic disaster recovery by a key holder. To amortize
network round-trips, many blocks are bundled into a single CAR~v1 file and
written with one \texttt{PUT}; serialization runs at roughly $3.8$\,GiB/s.

% ===========================================================================
\section{The Datom Model}
\label{sec:datom}

\kotoba's logical unit is the \emph{datom}, a 5-tuple in the Datomic tradition:
\begin{definition}[Datom]
A datom is a tuple $(E, A, V, T, \mathit{op})$ where $E\in\cid$ is the entity,
$A$ is an attribute (predicate) name, $V$ is a typed value (text, integer, float,
boolean, \cid, keyword, symbol, $\dots$), $T\in\cid$ is the transaction commit
that introduced the datom, and $\mathit{op}\in\{\textsf{assert},
\textsf{retract}\}$.
\end{definition}

Datoms are never updated in place; a logical change is a retraction plus an
assertion in a new transaction. Because $T$ is itself a content identifier and
transactions form an immutable, parent-linked \emph{CommitDag}, the entire
history is addressable and time-travel queries are simply scans restricted to the
datoms visible as of a given commit.

\paragraph{Write path.}
\texttt{QuadStore::transact} buffers asserted datoms, then on commit (i) builds
the four indices of \cref{sec:index} in parallel on large-stack threads,
(ii) rebuilds the four Prolly Trees, (iii) flushes the new blocks as a single
CAR~v1 bundle, and (iv) appends a commit node to the CommitDag, optionally
publishing the new head to IPNS for durability anchoring.

% ===========================================================================
\section{Indexing: The Four-Index Arrangement}
\label{sec:index}

Following Datomic, \kotoba\ maintains four covering indices so that any binding
pattern over $(E,A,V)$ resolves without a scan of the full datom set:
%
\begin{center}
\small
\begin{tabular}{@{}llp{0.42\linewidth}@{}}
\toprule
\textbf{Index} & \textbf{Order} & \textbf{Resolves} \\
\midrule
EAVT & $E\to A\to [V]$ & entity star, point $(E,A)$ lookup \\
AEVT & $A\to E\to [V]$ & predicate scan, reverse lookup \\
AVET & $A\to V\to [E]$ & value lookup ($A{=}a, V{=}v \Rightarrow E$) \\
VAET & $V\to A\to [E]$ & reverse references (ref-typed values) \\
\bottomrule
\end{tabular}
\end{center}
%
The hot arrangement holds these in tuned in-memory containers (hash maps for
point access, B-tree maps for ranged value/predicate scans); the cold Prolly Tree
mirrors the same four orders. A basic-graph-pattern (BGP) triple is routed to the
index that minimizes work given which of $E,A,V$ are bound---e.g.\ an unbound
subject with bound predicate and object goes to AVET, while a bound subject with
unbound predicate goes to EAVT. Multi-triple patterns route each triple
independently and intersect candidate entity sets before fetching full quads.
Transitive property paths (\texttt{$p$+}) are evaluated as bounded breadth-first
search over VAET/AEVT up to a configurable hop limit.

% ===========================================================================
\section{Incremental Datalog}
\label{sec:datalog}

The query engine evaluates Datalog programs by \emph{semi-naive} fixpoint
iteration over change deltas. Rather than re-deriving all facts at each
iteration, each rule is evaluated against the newly derived delta of the previous
round, and only genuinely new facts propagate; iteration halts when a round
produces nothing new. The canonical example is transitive closure,
\[
\textsf{reachable}(a,b) \;\leftarrow\;
  \textsf{edge}(a,b)
  \;\vee\;
  \bigl(\textsf{reachable}(a,x)\wedge\textsf{edge}(x,b)\bigr),
\]
which converges when no new \textsf{reachable} pair is derived. Because the
arrangement already maintains deltas per commit, incremental maintenance of
derived relations is a natural fit, and aggregate results can be memoized.

\paragraph{Materialized views.}
A query is cached under a \cid\ derived from the triple (head-commit \cid, graph
\cid, query text). A repeat of the same query against the same head is an $O(1)$
block fetch and decode; a new commit changes the key, transparently invalidating
the old view while keeping it addressable for historical queries.

\paragraph{SPARQL 1.1 surface.}
For interoperability, \kotoba\ exposes a SPARQL~1.1 frontend that parses queries
into algebra and translates them onto the datom indices. It supports
\texttt{SELECT}, \texttt{ASK}, \texttt{DESCRIBE}, \texttt{CONSTRUCT}, and
\texttt{UPDATE}, including filters, unions, optionals, sub-selects, property
paths, aggregation, and federated \texttt{SERVICE} calls routed by graph \cid.
\texttt{ASK} short-circuits on the first match and is effectively constant time.

% ===========================================================================
\section{Pregel BSP Computation}
\label{sec:pregel}

Iterative whole-graph algorithms (reachability, PageRank, clustering) run on a
Pregel bulk-synchronous-parallel engine. We map Pregel's abstractions directly
onto the datom graph: a vertex id is a subject \cid, vertex state is CBOR-encoded
facts, a message payload is a datom delta, and the per-vertex \texttt{compute}
step is an invocation of the incremental Datalog evaluator on the inbox. Each
superstep runs all active \texttt{compute} calls, waits at a barrier, delivers
messages, and repeats until every vertex has voted to halt and no messages remain.
A distributed runner shards vertices across nodes and exchanges inter-node
messages over libp2p GossipSub, sorting each inbox deterministically by source
\cid\ so that runs are reproducible.

% ===========================================================================
\section{Security Architecture}
\label{sec:security}

\subsection{CACAO capability delegation}
\label{sec:cacao}
Authorization is capability-based and chain-agnostic, following CAIP-74 (CACAO).
A request carries a delegation chain of CACAO objects, each with a header,
a payload (issuer and audience DIDs, issued-at and expiry timestamps, an
anti-replay nonce, and a list of resource URIs such as
\texttt{kotoba://graph/<cid>}), and a signature. For a depth-2 chain, the root
grant binds a grantor to an initial audience and the leaf invocation binds that
audience (now issuer) to the operator; verification checks both signatures,
enforces audience--issuer continuity, and enforces \emph{attenuation}---the leaf's
capabilities and graph scope must be a subset of the root's. Signatures may be
EdDSA (Ed25519, the native method), EIP-191/SIWE or EIP-1271 (EVM accounts), or
BIP-322 (Bitcoin), so credentials are portable across ecosystems.

Replay is prevented by a sharded nonce store, temporal validity by the
issued-at/expiry pair, and scope by requiring the requested graph \cid\ to appear
among the granted resources. The check itself is non-cryptographic and runs in
under a microsecond; signature verification dominates. Under load the gate sustains
the throughput reported in \cref{sec:eval} with full replay protection. Every
authorized private read may additionally emit an \emph{access receipt} datom
(graph, accessor DID, operation, purpose, timestamp), written in batches to an
audit subject that is itself versioned and replicated---providing an
X-Road-style accountability trail that is time-travel queryable.

\subsection{End-to-end encryption}
Agent-to-agent messages use the Signal protocol. Session establishment uses
Extended Triple Diffie-Hellman (X3DH): the initiator combines its identity and
ephemeral keys with the responder's identity, signed pre-key, and optional
one-time pre-key over X25519, and derives a shared secret with HKDF-SHA256 (with
the standard $\mathtt{0xFF}{\times}32$ prefix). Thereafter a Double Ratchet
provides forward secrecy and post-compromise security: a symmetric chain ratchet
derives a fresh AES-256-GCM message key per message, and a Diffie-Hellman ratchet
rotates root keys as new ratchet public keys arrive. Out-of-order and dropped
messages are tolerated by caching skipped message keys.

\subsection{Custody}
The block-sealing key (\cref{sec:store}) can be split among $N$ custodians using
Shamir secret sharing so that any $t$ of them can reconstruct it, decoupling
durability of ciphertext (on IPFS) from access to plaintext (gated by the
threshold).

% ===========================================================================
\section{WASM Runtime and Polyglot Compilation}
\label{sec:wasm}

User logic---validators, user-defined functions, and even small models---runs as
WebAssembly Component-Model components on a \texttt{wasmtime} host. The host
exposes a typed WIT world giving guests capability-scoped access to the database
(assert/retract/query/count datoms), the commit chain (link and signature
verification), authorization (capability checks), read-only EVM and Bitcoin
queries, and LLM inference. Execution is metered by a simple gas model (e.g.\
$10$ gas per assert, $100$ per query, $1000$ per inference, with a default ceiling
of $10^{7}$ gas), and each program is content-addressed by the \cid\ of its WASM
bytes and JIT-cached after first use. Guests may be written in any language that
targets the Component Model---Rust, Python, JavaScript, Go, C---and \kotoba\
ships a native \texttt{kotoba wasm} Kotoba/EDN-to-WASM compiler (implemented by
\texttt{kotoba-clj}) that parses EDN,
lowers to a typed AST, and emits WebAssembly in two passes, with a linear-memory
string/byte model and canonical-ABI host bindings sufficient to drive the Pregel
and LangGraph workflows end to end.

% ===========================================================================
\section{Evaluation}
\label{sec:eval}

We report representative measurements taken on an Apple~M4 laptop with release
builds. They are intended to characterize orders of magnitude, not to establish
competitive rankings; absolute numbers depend on graph size, concurrency, and the
cold-tier backend (in-memory, LAN Kubo, WAN Kubo, or same-AZ S3).

\begin{table}[h]
\centering
\small
\caption{Hot-path arrangement micro-benchmarks (in-memory block store).}
\label{tab:hot}
\begin{tabular}{@{}lr@{}}
\toprule
\textbf{Operation} & \textbf{Latency} \\
\midrule
EAVT point lookup            & $180$\,ns \\
AEVT reverse (B-tree range)  & $8$\,\textmu s \\
AVET value lookup            & $18$\,\textmu s \\
2-hop path                   & $748$\,ns \\
3-hop path                   & $1.16$\,\textmu s \\
JOIN (AVET$\times$AVET, 10K entities)   & $783$\,\textmu s \\
GROUP BY COUNT (50K entities)           & $47$\,ms \\
\bottomrule
\end{tabular}
\end{table}

\begin{table}[h]
\centering
\small
\caption{Macro-benchmarks: ingest and CACAO-gated query throughput.}
\label{tab:macro}
\begin{tabular}{@{}lr@{}}
\toprule
\textbf{Workload} & \textbf{Throughput} \\
\midrule
Batched ingest (sustained, 10 batches)        & $5{,}222$ entities/s \\
Batch ingest (1000 entities)                  & $3{,}981$ entities/s \\
\texttt{ASK} (constant-time, $c{=}32$)        & $12{,}753$ QPS \\
\texttt{SELECT} (result-bound, $c{=}8$)       & $3{,}777$ QPS \\
GROUP BY COUNT                                & $189$ QPS \\
2-triple JOIN                                 & $36.9$ QPS \\
HTTP \texttt{kg.sparql} \texttt{ASK}          & $2{,}586$ QPS \\
HTTP \texttt{kg.sparql} \texttt{SELECT} ($c{=}16$, 2000 ent.) & $398$ QPS \\
\bottomrule
\end{tabular}
\end{table}

\Cref{tab:hot} shows that hot lookups are sub-microsecond and that multi-hop
traversals stay in the single-digit-microsecond range, reflecting the covering
indices. \Cref{tab:macro} shows that constant-time \texttt{ASK} queries sustain
over $12$K QPS even with full CACAO verification and replay protection, while
result-bound \texttt{SELECT}s and aggregates scale with result size as expected.
Cold-path lookups add the block-store round-trip: in our measurements a cold
prefix scan cost about $3.3$\,ms against a LAN Kubo node, $6.0$\,ms against
same-AZ S3, and $175$\,ms against a WAN Kubo node, dominated by tree depth
($\approx 3$ for thousands of entries) times the per-block latency.

% ===========================================================================
\section{Related Work}
\label{sec:related}

\kotoba\ draws on several lineages. Its immutable, time-traveling datom model and
four-index design follow Datomic~\cite{datomic}. Content addressing, IPLD,
DAG-CBOR, and the CAR container format come from IPFS~\cite{ipfs,ipld} and
libp2p~\cite{libp2p}. The Prolly Tree is a probabilistic B-tree popularized by
Dolt/Noms for history-independent, diffable storage~\cite{prollytree}.
Incremental Datalog~\cite{datalog} and arrangement-based delta maintenance are in
the tradition of differential dataflow~\cite{differential-dataflow}. The SPARQL
surface targets the W3C RDF query standard~\cite{sparql}.
Pregel~\cite{pregel} provides the BSP model for whole-graph computation. CACAO
(CAIP-74)~\cite{cacao}, EIP-4361/SIWE~\cite{siwe}, and DIDs~\cite{did} provide
chain-agnostic object capabilities, and the AT Protocol~\cite{atproto} informs the
XRPC service and firehose surfaces. The Signal protocol (X3DH~\cite{signal-x3dh}
and Double Ratchet~\cite{signal-doubleratchet}) provides end-to-end message
security; Shamir secret sharing~\cite{shamir} underlies threshold custody; the
WebAssembly Component Model with WIT~\cite{wasm-component} provides the polyglot
sandbox, into which the Kotoba/EDN compiler emits. Kotoba/EDN retains
Lisp/EDN-shaped syntax and data notation rooted in the Clojure
ecosystem~\cite{clojure}. \kotoba's contribution is not any one of these
mechanisms but their coherent integration behind a single content-addressed
datom substrate.

% ===========================================================================
\section{Limitations and Roadmap}
\label{sec:roadmap}

The current deployment is effectively single-writer per graph: conflicting
commits are rejected by compare-and-swap rather than merged. Our roadmap
addresses this along five orthogonal axes:
%
\begin{enumerate}[leftmargin=1.4em,itemsep=2pt]
  \item \textbf{Authority}---from owner-rooted CACAO to delegatable grants with
    per-graph WASM validation hooks.
  \item \textbf{Ordering}---from wall-clock timestamps to hybrid logical clocks
    on commits, with causal cross-references.
  \item \textbf{Availability}---from a de-facto owner node to declared
    per-graph replication policies enforced by a DHT neighborhood.
  \item \textbf{Conflict}---from CAS rejection to merge-on-conflict via a
    multi-parent commit DAG and Merkle-CRDT semantics (last-write-wins / OR-set).
  \item \textbf{Transport}---from live-only non-datomic streams to fully
    cursor-replayable pointer-chains for all topics.
\end{enumerate}
%
The storage-confidentiality roadmap is similarly staged: encryption at rest and
access receipts are implemented, while per-DID signed journals anchored to an L2
and threshold custodian rotation are planned.

% ===========================================================================
\section{Conclusion}
\label{sec:conclusion}

\kotoba\ shows that a single content-addressed datom substrate can carry the
storage, query, computation, authorization, encryption, and execution concerns of
a decentralized AI-agent system. By making facts, indices, transactions, and
programs all addressable IPFS objects, the database becomes a Merkle-CRDT whose
identities are coordination-free, while a hot arrangement keeps interactive
queries fast and CACAO, Signal, and WASM provide portable security and safe
extensibility. Our evaluation indicates the design is practical on commodity
hardware, and our five-axis roadmap charts the remaining work toward full
multi-writer decentralization.

% ---------------------------------------------------------------------------
\paragraph{Availability.}
Source, documentation, and interactive explainers are available at
\url{https://com-junkawasaki.github.io/kotoba/}.

\bibliographystyle{plain}
\bibliography{references}

\end{document}
